\begin{theorem}[The cusp replacement and its relative homotopy class]\hypertarget{tag-02PD}{}\label{mapsI:src:A.2.1}
\StatementRecord{mapsI:src:A.2.1}{proved}{02PD}{\SeriesRef{mapsI:quotient-eq-5}}

For every admissible member, there is a smooth map
$e_{\rm q}:X\to Q$ over $B$ whose restriction outside an arbitrarily
small cusp neighbourhood is the literal circle quotient.
For fixed quotient gluing, its homotopy class is independent
of the cusp replacement, provided the replacement maps the
cusp neighbourhood into its designated solid torus and is
kept fixed on its outer collar.  In particular there is a
well-defined class for the original member $X\simeq S^6$.
\end{theorem}



\begin{proof}
\begin{enumerate}[label=\arabic*.,leftmargin=*]
\item On the cusp boundary use the product coordinates
$(q,\theta_\gamma)$ after the specified collar shear.
By \eqref{mapsI:quotient-eq-5}, the boundary circle class extends to $N_0$.
Since $S^1=K(\mathbb Z,1)$, it has a continuous representative
$\theta':N_0\to S^1$.  Its restriction to the boundary is
homotopic to the given one: the two maps induce the same
homomorphism on $\pi_1$ of the boundary mapping torus.
The homotopy-extension property for a manifold with boundary
changes $\theta'$ on a collar so that it agrees exactly with
$\theta_\gamma$ there.  Relative smooth approximation gives
a smooth representative, still equal on a smaller collar.
Then $(f,\theta')$ is the required map into $D^2\times S^1$,
and agrees with the finite-filling maps on the overlap.
Both components of $(f,\theta')$ are smooth on $N_0$, so the
product map is smooth at the critical points of $f$ as well.

\item The inclusion $\partial N_0\to N_0$ surjects on $\pi_1$,
as \eqref{mapsI:quotient-eq-5} explicitly shows.  It follows that
$H^1(N_0;\mathbb Z)\to H^1(\partial N_0;\mathbb Z)$
is injective.  Since both spaces are connected, the long
exact sequence gives $H^1(N_0,\partial N_0;\mathbb Z)=0$.
Two circle replacements agreeing on the boundary therefore
have quotient whose relative $H^1$ class is zero, hence
have a relative homotopy.  Keeping the base coordinate
fixed gives the required homotopy of maps into the solid
torus.  Different choices of the collar size are compared
on a common smaller collar.
\end{enumerate}
\end{proof}